% !PRINTED 399
\weberpage{399}
\begin{center}
\S.\,128.

\gesperrt{Derivation of all solutions of the Pell equation
from the least positive one.}
\end{center}

If $t$, $u$ is any solution of the Pell equation, then
the expression
\[
\frac{t^2 - D u^2}{4},
\]
which has the value $\pm 1$, is the norm of the two conjugate numbers
\[
\frac{t + u\sqrt{D}}{2}, \quad \frac{t - u\sqrt{D}}{2}.
\]
We shall call them the \gesperrt{units} belonging to $\sqrt{D}$.
If, then, $(t_1, u_1)$, $(t_2, u_2)$ are any two (positive or negative)
solutions of the Pell equation
\[
t^2 - D u^2 = \pm 4,
\]
then the two numbers
\[
\Theta_1 = \frac{t_1 + u_1 \sqrt{D}}{2}, \quad \Theta_2 = \frac{t_2 + u_2 \sqrt{D}}{2}
\]
have the norm $\pm 1$, and the same therefore holds also of their product
\[
\Theta_3 = \Theta_1 \Theta_2 = \frac{t_3 + u_3 \sqrt{D}}{2}.
\]
Here
\[
t_3 = \frac{t_1 t_2 + u_1 u_2 D}{2}, \quad u_3 = \frac{t_1 u_2 + t_2 u_1}{2},
\]
and $t_3$ and $u_3$ are, as one sees at once, integers.

For if $D$ is even, then $t_1$, $t_2$ are even, and if $D$
is odd, then $t_1$, $u_1$ are either both even or both
odd, and likewise $t_2$, $u_2$. Thus $\Theta_3$ is likewise a unit
belonging to $\sqrt{D}$, and it follows \gesperrt{that the product of two}
(and consequently also of several) \gesperrt{units is again a unit.}

If $\Theta$ is an arbitrary unit, then $\pm \Theta^{-1}$ is the unit conjugate
to $\Theta$. Among the units different from $\pm 1$
there always belong four together,
\[
\pm \Theta, \quad \pm \Theta^{-1},
\]
among which \gesperrt{one} is positive and greater than 1. This one corresponds
%CONT>>
